\scanpage{10}
\mspage{P 50}
profonde à installer une pièce après l'autre, et à voir peu à peu s'instaurer, parmi ces salles, ces chambres et ces réduits l'ordre harmonieux de la maison vivante - belle, accueillante, bonne pour y vivre. Car l a \quad m a i s o n, avant tout et secrètement en chacun de nous, c'est aussi l a \quad m è r e - ce qui nous entoure et nous abrite, à la fois refuge et réconfort ; et peut-être (plus profondément encore, et alors même que nous serions en train de la construire de toutes pièces) c'est cela aussi dont nous sommes nous-mêmes issus, ce qui nous a abrité et nourri, en ces temps à jamais oubliés d'avant notre naissance... C'est aussi l e \quad G i r o n .

Et l'image apparue spontanément tantôt, pour aller au delà de l'appellation prestigieuse de ``pionnier'', et pour cerner la réalité plus cachée qu'elle recouvrait, était elle aussi dépouillée de tout accent ``héroïque''. Là encore, c'était l'image archétype du maternel qui est apparue - celle de la ``matrice'' nourricière et de ses informes et obscurs labeurs...

Ces deux pulsions qui m'apparaissaient comme ``de nature différente'' sont finalement plus proches que je ne l'aurais pensé. L'une et l'autre sont dans la nature d'une `` p u l s i o n \quad d e \quad c o n t a c t '' , nous portant à la rencontre de `` l a \quad M è r e '' : de Celle qui incarne e t ce qui est proche, ``connu'', e t ce qui est ``inconnu''. M'abandonner à l'une ou l'autre pulsion, c'est `` r e t r o u v e r \quad l a \quad M è r e ''. C'est renouveler le contact à la fois au p r o c h e , au ``plus ou moins connu'', et au `` l o i n t a i n '', à ce qui est ``inconnu'' mais en même temps pressenti, sur le point de se faire connaître.

La différence ici est de tonalité, de dosage, non de nature. Quand je ``bâtis des maisons'', c'est le ``connu'' qui domine, et quand ``j'explore'', c'est l'inconnu. Ces deux ``modes'' de découverte, ou pour mieux dire, ces deux aspects d'un même processus ou d'un même travail, sont indissolublement liés. Ils sont essentiels l'un et l'autre, et complémentaires. Dans mon travail mathématique, je discerne un mouvement de va-et-vient constant entre ces deux modes d'approche, ou plutôt, entre les moments (ou les périodes) où l'un prédomine, et ceux où prédomine l'autre (*). Mais il est clair aussi qu'en chaque moment, et l'un et

\noindent\rule{2.1cm}{0.4pt}

(*) Ce que je dis ici sur le travail mathématique est vrai également pour le travail de ``méditation'' (dont il sera question un peu partout dans Récoltes et Semailles). Il n'y a guère de doute pour moi que c'est là une chose qui apparaît dans tout travail de découverte, y compris dans celui de l'artiste (écrivain ou poète, disons). Les deux ``versants'' que je décris ici peuvent être vus également comme étant, l'un celui de l ' e x p r e s s i o n et de ses exigences ``techniques'', l'autre celui de la r é c e p t i o n (de perceptions et d'impressions de toutes sortes), devenant i n s p i r a t i o n par l'effet d'une attention
% archive sheet number: 9

